\section{Evolução do projeto}


\begin{frame}

    \frametitle{Evolução do projeto}
    
    \begin{itemize}
    	
    	\item Definidas na última apresentação:
		\begin{itemize}
			\item Reconstruir e ampliar o dataset.	
			\item Realizar novos treinamentos.	
			\item Otimizar hiperparâmetros com Optuna.	
			\item Realizar novas análises com Grad-CAM.
			\item Avaliar desempenho, generalização e interpretabilidade.			
		\end{itemize}
		
		 \item Extras:
		\begin{itemize}
			\item Exportação do modelo e sua implantação em um aplicativo Android.
		\end{itemize}
	\end{itemize}
\end{frame}

\begin{frame}
	
	\frametitle{Construção do dataset integrado}
	
	Foram integradas quatro bases públicas de imagens de resíduos, que apresentavam diferenças nos rótulos, quantidade de classes, resolução e características das imagens.

	\textbf{Etapas realizadas:}
	\begin{itemize}
		\item Harmonização dos rótulos entre as bases.
		\item Curadoria semântica utilizando CLIP.
		\item Identificação e remoção de imagens duplicadas.
		\item Divisão dos dados por grupos de imagens, evitando vazamento entre treino, validação e teste.
	\end{itemize}
	
\end{frame}

\begin{frame}
	
	\frametitle{Construção do dataset integrado}
	
	Resultado da integração:
	\begin{itemize}
		\item 4 bases públicas.
		\item 35.603 imagens indexadas.
		\item 3.121 imagens redundantes.
		\item 15.707 imagens no dataset combinado (fiquei em duvida dessa info).
	\end{itemize} 
	
\end{frame}